% Paper skeleton for paperreview.ai (venue ICLR), due 11 Dec 2026.
% With iclr2027_conference.sty in this folder the ICLR style is used; without it
% the same text compiles as a plain article, so the skeleton can be checked now.
\documentclass{article}
% Course copy: author name shown, course header. Set \coursecopyfalse for an
% anonymous ICLR submission (line numbers, "Under review" header).
\newif\ifcoursecopy \coursecopytrue
\newif\ificlrstyle
\IfFileExists{iclr2027_conference.sty}{\iclrstyletrue}{\iclrstylefalse}
\ificlrstyle
  \usepackage{iclr2027_conference,times}
  \ifcoursecopy\iclrfinalcopy\fi
  \newcommand{\mybibstyle}{iclr2027_conference}
\else
  \usepackage{times}\usepackage[margin=1in]{geometry}\usepackage[round]{natbib}
  \newcommand{\mybibstyle}{plainnat}
\fi
\usepackage{amsmath,amssymb}
\usepackage{graphicx}
\usepackage{booktabs}
\usepackage{url}
\usepackage[hidelinks]{hyperref}
\usepackage{xcolor}
\usepackage{tikz}
\usetikzlibrary{positioning,arrows.meta,fit,calc}
\newcommand{\todo}[1]{\textcolor{red}{[TODO: #1]}}

\title{Wait-Aware Steering of MCP Tool Calls\\across 5G MEC Sites}
\author{Bui Duong Kham (1102549018)\\
Graduate School of Information and Communication Technology, Soongsil University}

\begin{document}
\maketitle
\ificlrstyle\ifcoursecopy\lhead{Term project, Machine Learning and AI, Fall 2026}\fi\fi

\begin{abstract}
An AI agent on a phone finishes one task with a chain of tool calls, and operators start
to host the tool servers (MCP servers) at several edge sites. The mobile network can pick
a site for every new call, but today it picks the nearest site or the one with the fewest
calls in flight. Neither sees that one SQL tool needs 5\,ms for a lookup and several
seconds for a scan, so a lookup can wait behind two scans. We let each tool replica
predict the run time of every call from its arguments, compute how long a new call would
wait for a free worker, and report this wait to a controller that steers new connections
at the 5G user-plane gateway (UPF). On a free5GC testbed with 30 emulated phones and
replicas at three edge sites, this lowers the 95th-percentile call latency by 16.1\%
compared with least-load steering (95\% CI 13.2--19.5\%, 10 paired runs), and by 20\%
for short calls. A baseline with the same wait formula but a per-tool mean run time gains
nothing, so the gain comes from reading the arguments. \todo{one sentence on the load
$\times$ skew crossover: where steering pays off and where the nearest site is enough.}
\end{abstract}

\section{Introduction}
\todo{1 page. (1) Agents call tools in chains: 10--30 dependent calls per task, so the
slowest call decides the task time; with 20 calls, $1-0.95^{20}\approx 64\%$ of tasks
meet at least one call above the per-call p95, so we report p95 rather than the mean.
(2) Run times of one tool spread over three orders of magnitude (TraceLab: median 0.3\,s,
p99 3.6\,min). (3) The network can steer, but sees only addresses: call arguments are
encrypted. (4) Our idea: the replica reads its own queue and reports a wait; the network
uses the wait. (5) Contributions, with the measured numbers.}

Contributions:
\begin{itemize}
  \item A wait estimator at each tool replica: it predicts the run time of each queued
  call from its arguments and replays the queue over the workers to give the time until
  a worker is free for a new call (Section~\ref{sec:method}).
  \item A steering controller at the UPF that splits new connections over the sites by
  that wait plus the network round trip, with the same code for every baseline.
  \item Measurements on a real 5G core with real MCP traffic: the main comparison, an
  ablation with exact run times and with argument-free means, the per-call connection
  cost, and a load $\times$ skew sweep that shows where steering pays off
  (Section~\ref{sec:results}).
\end{itemize}

\begin{figure}[t]
\centering
\begin{tikzpicture}[
  box/.style={draw, rounded corners=2pt, align=center, font=\small, minimum height=8mm},
  site/.style={draw, rounded corners=2pt, align=left, font=\small, text width=33mm, minimum height=13mm},
  >={Stealth[length=2mm]}, node distance=6mm and 9mm]
  \node[box, text width=21mm] (ue) {UE\\agent + MCP client};
  \node[box, right=of ue] (gnb) {gNB};
  \node[box, right=of gnb, text width=24mm] (upf) {UPF\\steering map:\\anycast $\to$ site};
  \node[site, right=16mm of upf, yshift=17mm] (m1) {MEC-1, RTT 1.7\,ms\\replica: FIFO queue,\\2 workers, wait estimator};
  \node[site, below=4mm of m1] (m2) {MEC-2, RTT 11.9\,ms\\replica (same code)};
  \node[site, below=4mm of m2] (m3) {MEC-3, RTT 29.7\,ms\\replica (same code)};
  \node[box, below=14mm of upf, text width=24mm] (ctl) {controller\\every 100\,ms};
  \draw[->] (ue) -- node[above, font=\scriptsize] {call} (gnb);
  \draw[->] (gnb) -- (upf);
  \draw[->] (upf.east) -- (m1.west);
  \draw[->] (upf.east) -- (m2.west);
  \draw[->] (upf.east) -- (m3.west);
  \node[draw, dotted, fit=(m1)(m2)(m3), inner sep=2mm] (sites) {};
  \draw[->, dashed] (sites.south) |- node[pos=0.75, below=2pt, font=\scriptsize] {wait report of every site} (ctl.east);
  \draw[->, dashed] (ctl) -- node[left, font=\scriptsize, align=right] {new split\\of new\\connections} (upf);
\end{tikzpicture}
\caption{Setup. Each MCP server has one anycast address. The UPF sends every new
connection (one call) to one site by a split map. Each replica reads the arguments of
its queued calls, predicts their run times and reports the wait for a new call; the
controller turns the reports into the split of the next interval.}
\label{fig:setup}
\end{figure}

\section{Background and related work}
\todo{0.5--1 page. Terms in plain words: UE, UPF, MEC site, anycast, CATS, MCP over
Streamable HTTP (stateless since the 2026-07-28 spec, so a per-call decision does not
split a session).}

\textbf{Steering by load.} IETF CATS \citep{cats_rfc,cats_metrics} picks an anycast copy by
network distance and a server load metric: a request count, a utilization or an average.
DONAR \citep{donar}, Taiji \citep{taiji}, Envoy load-aware locality \citep{envoy} and
\citet{fan2018} split load across sites. We use this kind of split; our change is the
signal.

\textbf{Tool calls of agents.} DOME \citep{dome} predicts the delay of each MCP call and
spreads calls over tool instances, but in an application dispatcher inside one LAN.
EdgeAgent/WasmMCP \citep{edgeagent} places each call on device, edge or cloud from a
profile, without load. 3GPP TR 26.870 \citep{tr26870} runs MCP sessions under emulated 5G
conditions without a UPF or steering. \todo{agent benchmarks: ReAct \citep{react},
$\tau$-bench \citep{taubench}; they treat a call as instant.}

\textbf{Gap.} None of these ranks sites by a wait predicted from the arguments of the
queued calls, and none measures MCP traffic through a real 5G UPF to several sites.

\section{Method}
\label{sec:method}
\textbf{Run-time prediction.} For each call the replica builds features from the arguments
only: the SQL query type, the tables and joins, the row estimate from \texttt{EXPLAIN},
\texttt{LIMIT}, and the tool. A gradient-boosted tree model \citep{lightgbm} predicts the
log run time. An online factor (moving average of measured / predicted run time) corrects
the prediction for the current speed of the site. \todo{held-out $R^2$ table: by task 0.88
on log run time, 0.54 in seconds; by template 0.62 / 0.10.}

\textbf{Wait for a new call.} The replica knows the predicted remaining run time of each
running call and the predicted run time of each queued call. It replays the FIFO queue on
its $w$ workers (a heap of worker free times) and reports the time until a worker is free
for a new call. A simpler formula, remaining work divided by $w$, is wrong for small $w$:
one long call next to an idle worker gives no wait at all. \todo{one equation.}

\textbf{Split.} Every 100\,ms the controller reads the latest report of each site and
ignores reports older than 300\,ms. It estimates how many calls arrive in the next 200\,ms
and assigns them to sites by water-filling on (reported wait + RTT), so that one stale view
does not send a herd of calls to one site. The split becomes a map from a random bucket
0--99 to a site in the UPF; every new connection draws one bucket.

\textbf{Baselines.} B0 sends every call to the nearest site (plain anycast). B1 uses the
number of calls in flight (CATS-style). B2 uses the same wait formula as our method with
the mean run time of each tool, without the arguments. All policies use the same split,
period and horizon; only the signal changes. The oracle uses the measured run time of
each call (an upper bound for the predictor).

\section{Testbed and workload}
\label{sec:setup}
\todo{0.5 page. 5G core free5GC v4.3.0 with one UPF, UERANSIM gNB and 30 UEs; three MEC
sites with RTT 1.7 / 11.9 / 29.7\,ms (Table: RTT and one MCP call on a new connection);
replicas: SQL tool on TPC-H SF10 (DuckDB, 2 workers) and a 5\,ms key-value tool. Workload:
5,119 calls recorded from a Qwen3-8B agent on 50 analytics tasks (10 templates $\times$ 5
parameter sets), replayed by 90 clients (3 agents per UE) at a set share of the measured
capacity; local users take a set share of MEC-1. Say which traffic is live and which is
replayed. Runs: 300\,s, first 60\,s dropped, the same seed offers the same calls to every
policy.}

\section{Results}
\label{sec:results}

\begin{table}[t]
\centering
\caption{Main setting: offered work 70\% of capacity, local users take 30\% of MEC-1.
Means over runs; gains are paired by seed ($1-\text{policy}/\text{B1}$) with a bootstrap
95\% CI over seeds. Short calls: run time below 50\,ms.}
\label{tab:main}
\small
\begin{tabular}{lrrrrl}
\toprule
Policy & Runs & p50 (ms) & p95 (ms) & p95 short (ms) & p95 gain vs B1 \\
\midrule
B0 nearest & 3 & 2974 & 10158 & 9970 & $-312\%$ \\
B1 fewest in flight & 10 & 71 & 2483 & 1897 & --- \\
B2 per-tool mean & 10 & 72 & 2503 & 1905 & $-1.2\%$ [$-4.7$, 2.1] \\
Ours & 10 & 73 & \textbf{2085} & \textbf{1503} & \textbf{16.1\%} [13.2, 19.5] \\
Oracle run times & 10 & 75 & 2079 & 1466 & 16.5\% [13.4, 19.5] \\
\bottomrule
\end{tabular}
\end{table}

\textbf{Main comparison.} Table~\ref{tab:main}. Our method lowers p95 by 16.1\% compared
with B1 and wins in all 10 runs. The gain comes from short calls (p95 $-20.2\%$): a lookup
no longer lands behind a scan that B1 cannot see. The median is 2.5\% worse, because some
calls travel to a farther site. \todo{p95 without any queue wait (latency minus wait) as
the floor.}

\textbf{What the gain comes from.} B2 uses the same wait formula without the arguments and
does not beat B1, so the gain comes from reading the arguments. With exact run times the
gain is 16.5\%, almost the same: here a better predictor would add little.

\textbf{When steering pays off.} \todo{Figure: load 10--90\% $\times$ local-user share
0/30/60\% at MEC-1, p50 and p95 per policy, crossover with B0 marked. Seed-1 reading:
median crossover at about 23\% load with no local users, about 14\% with 30\%, below 10\%
with 60\%; at 90\% load the policies converge.}

\begin{table}[t]
\centering
\caption{Cost of one new connection per call: 200 key-value calls per cell from one UE
through the UPF (p50).}
\label{tab:conn}
\small
\begin{tabular}{lrrr}
\toprule
Site & New connection (ms) & Kept-alive connection (ms) & Setup cost (ms) \\
\midrule
MEC-1 & 8.4 & 8.1 & 0.3 \\
MEC-2 & 29.8 & 19.8 & 10.0 \\
MEC-3 & 76.5 & 45.5 & 31.0 \\
\bottomrule
\end{tabular}
\end{table}

\textbf{Connection cost.} Table~\ref{tab:conn}. Steering per call needs one new connection
per call, which costs one round trip. It is small next to the waits above, but not next to
a 5\,ms lookup at low load.

\textbf{Failure cases.} \todo{from the sweep: at low load the method and B1 tie on p95 and
sending a short call 30\,ms away is slower than waiting; at 90\% load every site is busy
and moving a call does not help; cells where the method loses to B1.}

\section{Limitations}
\todo{A gateway at the MEC that reads requests (L7) could also steer; replayed rather than
live agents create the load; one host runs all sites, so long scans slow down when all
sites are busy; one SQL tool and one key-value tool.}

\section{Conclusion}
\todo{3 sentences.}

\section*{Use of AI}
AI assistants were used to search for papers by keyword, to review the feasibility of the
plan, and to suggest the tech stack and code.

\bibliographystyle{\mybibstyle}
\bibliography{refs}
\end{document}
